% ==================================================================
% id:       app-020
% topic:    שימושי הנגזרת
% subtopic: פולינום טיילור
% solution: full
% ==================================================================
רוצים לחשב את
\(\sin (0.3)\) בדיוק של 5 ספרות לפחות. קבעו את המעלה של פולינום מקלורן שאותו יש לקחת,
ומצאו את הקירוב הרצוי.
\hint{השתמשו בשארית לגרנז' וחסמו את הנגזרות של $\sin$ בערך מוחלט. שימו לב שבפיתוח של $\sin$ מופיעות רק חזקות אי-זוגיות.}
\begin{solbox}
נזכר בפולינום מקלורן של
$\sin(x)$
מסדר זוגי:
\[
.P_{2n+2}(x)=x-\frac{x^3}{3!}+\frac{x^5}{5!}-\frac{x^7}{7!}+\cdots+(-1)^n\frac{x^{2n+1}}{(2n+1)!}
\]
מעוניינים במציאת
$n$
כך ש-
$$.|R_{2n+2}(x)|<10^{-5}$$
נזכר כי שארית לגרנז' נתונה בצורה
\begin{align*}
|f(x)-P_{2n+2}(x)|&=\ |R_{2n+2}(0.3)|\\
&=\ \left|(-1)^{n+1}\frac{\sin^{(2n+3)}(c)}{(2n+3)!}\cdot0.3^{2n+3}\right|\\
&\leq\ \frac{0.3^{2n+3}}{(2n+3)!},
\end{align*}
עבור
$0<c<0.3$
כלשהו.
עלינו למצוא פתרון ב-
$n$
לאי-שיוויון
$$.\frac{0.3^{2n+3}}{(2n+3)!}<\frac{1}{100000}$$
עתה נציב
ונבדוק איזה
$n$
יקיים את הנדרש.
בבירור
$n=1$
לא מתאים, הרי ש-
$$,|R_4(0.3)|\leq\frac{0.3^5}{5!}=\frac{2025}{10^8}$$
טוב, כאן אנו רואים כי המספר שקיבלנו -
$\frac{2025}{10^8}$
גדול מ-
$10^{-5}$
(שימו לב, כי זה לא אומר שהשארית גדולה מ-
$10^{-5}$!).
לא קשה להשתכנע כי
$n=2$
מהווה פתרון, אכן
\[
|R_6(0.3)|\le \frac{0.3^7}{7!}\approx 4.34\cdot 10^{-8}<10^{-5}.
\]
לכן צריך פולינום מקלורן \textbf{מעלה 6}:
$$\sin0.3\approx P_6(0.3)=0.3-\frac{0.3^3}{6}+\frac{0.3^5}{120}\approx0.2955206$$
\end{solbox}
